\chapter{Coprocessor Instructions}
\label{ch:coprocessor}

\section{Overview}

The SIRC-1 CPU includes a coprocessor interface that allows delegation of certain operations to specialized
coprocessors. The primary coprocessor is the Exception Unit (Coprocessor 1), which handles exceptions, interrupts, and
privileged operations.

Coprocessor calls are ordinary conditional instructions. If the instruction's condition is true, the processing unit
computes a command operand and writes it to the 16-bit pending coprocessor-command latch during write-back. The
selected coprocessor then interprets the command value. If the selected coprocessor is not present in the CPU model or
does not implement the requested operation, the exception unit raises an invalid-opcode fault so the operation can be
emulated or rejected by software.

\section{Legal Forms}

\begin{table}[H]
	\centering
	\small
	\begin{tabular}{llll}
		\toprule
		\textbf{Instruction} & \textbf{Syntax}       & \textbf{Encoding} & \textbf{Notes}                           \\
		\midrule
		COPI                 & \texttt{COPI \#imm16} & 0x0F              & 16-bit coprocessor command operand       \\
		COPR                 & \texttt{COPR rS}      & 0x3F              & Command operand is read from \texttt{rS} \\
		\bottomrule
	\end{tabular}
	\caption{Legal Coprocessor Instruction Forms}
	\label{tab:coprocessor-legal-forms}
\end{table}

\noindent The CPU encoding reserves the register field of the immediate format (opcode 0x0F) and the R1, R2, and shift
fields of the register format (opcode 0x3F); the immediate format has no shift fields. The programmer-facing
\mnemonic{COPI} and \mnemonic{COPR} syntax hides those fields; assemblers emit canonical values and coprocessors must
not interpret them as destination-register writes. Opcode 0x2F is Undocumented, in the same class as opcodes
\opcode{08}, \opcode{09}, \opcode{0B}, and \opcode{0D} (see Appendix~\ref{appendix:undocumented}); it has no public
assembly syntax.

\section{Common Semantics}

\begin{table}[H]
	\centering
	\begin{tabular}{lp{9cm}}
		\toprule
		\textbf{Topic}     & \textbf{Definition}                                                                                                                                                                                                                                                                                                                                                                                         \\
		\midrule
		Condition false    & The instruction is skipped and has no side effects. The pending coprocessor command, registers, memory, address registers, and status flags are preserved. The program counter still advances to the following instruction.                                                                                                                                                                                 \\
		Command operand    & \mnemonic{COPI} uses the 16-bit immediate value. \mnemonic{COPR} uses its source register value as the command operand.                                                                                                                                                                                                                                                                                     \\
		Command format     & For standard coprocessor commands, bits 15--12 select the coprocessor ID, bits 11--8 select the coprocessor operation, and bits 7--0 are coprocessor-defined operand bits.                                                                                                                                                                                                                                  \\
		Write-back         & The processing unit writes the command operand to the pending coprocessor-command latch. It does not write general-purpose registers, address registers, or data memory as part of the call. Hidden encoding fields are ignored.                                                                                                                                                                            \\
		Status flags       & Preserved by the processing-unit coprocessor-call write-back path. A coprocessor operation may define its own later CPU-state effects; exception-unit effects are documented in Chapter~\ref{ch:exceptions}.                                                                                                                                                                                                \\
		Privilege          & In protected mode, coprocessor-operation nibbles 0x0--0x7 are user-callable and 0x8--0xF are supervisor-only. Privilege checks only apply when the instruction condition is true. For \mnemonic{COPR} the privilege decision is made on the run-time value of \reg{rS} sampled during decode, not on any encoded field, so whether a given \mnemonic{COPR} instruction faults depends on register contents. \\
		Timing             & The coprocessor-call instruction itself uses the normal 6 execution phases and performs no data-memory bus access. Dispatch to coprocessor 1 then occupies one further complete 6-cycle slot beginning at the next phase 0, so a coprocessor-1 operation completes 12 cycles after the call's own fetch began.                                                                                              \\
		Exceptions         & Supervisor-only operations can raise privilege-violation faults; missing coprocessors and unimplemented operations can raise invalid-opcode faults. See the notes below.                                                                                                                                                                                                                                    \\
		Fault side effects & If a fault is raised while dispatching the command, software must not depend on the contents of the pending command latch after the fault. A pending coprocessor command is discarded if a fault or hardware exception is dispatched before the coprocessor's dispatch slot begins; the command is not re-issued after the handler returns.                                                                 \\
		\bottomrule
	\end{tabular}
	\caption{Common Coprocessor Instruction Semantics}
	\label{tab:coprocessor-common-semantics}
\end{table}

\noindent Coprocessor-call instructions do not perform data-memory access and do not raise data bus, data
bus-protection, data alignment, or effective-address segment-overflow faults. A program-counter wrap with
\texttt{SR.A} set still raises a segment-overflow fault at the following instruction fetch, as it does for every
instruction. Instruction fetch can still raise the normal fetch faults before the call executes. If trace mode was
enabled when the instruction began, an instruction-trace fault is raised after the processing-unit call commits,
unless an earlier fault prevents the instruction from completing. In protected mode, supervisor-only coprocessor
operation nibbles are rejected before the command is dispatched. Missing coprocessors and unimplemented coprocessor
operations are detected after the command has been latched for dispatch and raise invalid-opcode faults through the
exception unit.

\begin{instructionbox}{COPI/COPR}{Coprocessor Call}
	\textbf{Applicability:} SIRC-1, all revisions

	\textbf{Opcodes:} 0x0F (Immediate), 0x3F (Register). Opcode 0x2F is Undocumented. See Table~\ref{tab:coprocessor-legal-forms} for the full family of legal forms.
	\begin{center}
		\begin{bytefield}[bitwidth=1.1em]{32}
			\bitheader{0-31} \\
			\bitbox{6}{\opcode{0F}} &
			\bitbox{4}{0x0} &
			\bitbox{16}{Command} &
			\bitbox{2}{00} &
			\bitbox{4}{Cond}
		\end{bytefield}
	\end{center}
	\begin{center}
		\begin{bytefield}[bitwidth=1.1em]{32}
			\bitheader{0-31} \\
			\bitbox{6}{\opcode{3F}} &
			\bitbox{4}{0x0} &
			\bitbox{4}{0x0} &
			\bitbox{4}{rS} &
			\bitbox{1}{0} &
			\bitbox{3}{000} &
			\bitbox{4}{0x0} &
			\bitbox{2}{00} &
			\bitbox{4}{Cond}
		\end{bytefield}
	\end{center}

	\textbf{Instruction Fields:}
	\begin{itemize}
		\item \textbf{Register (bits 25--22, Immediate format only):} reserved; the assembler emits \texttt{0x0} and
		      coprocessors must not interpret it as a destination register.
		\item \textbf{Register 1, Register 2 (Register format, bits 25--18):} reserved; the assembler emits
		      \texttt{0x0} for both.
		\item \textbf{Register 3 (Register format, bits 17--14):} \texttt{rS}, any of the 16 registers,
		      \texttt{0x0}--\texttt{0xF}; its value becomes the 16-bit command operand.
		\item \textbf{Immediate (bits 21--6, Immediate format only):} 16-bit command operand for \mnemonic{COPI}.
		\item \textbf{Shift fields (Register format, bits 13--6):} reserved; the assembler emits \texttt{0} and
		      coprocessors must not interpret them.
		\item \textbf{AF (bits 5--4):} reserved; the assembler emits \texttt{00} and coprocessors must not interpret
		      it as a status-update selector.
		\item \textbf{Condition field (bits 3--0):} any of the 16 condition-code encodings in
		      Table~\ref{tab:condition-code-encoding}.
	\end{itemize}

	\textbf{Syntax:}
	\begin{lstlisting}
COPI #opcode                  ; Call coprocessor with 16-bit command
COPR rS                       ; Call coprocessor with command in rS
\end{lstlisting}

	\textbf{Operands:} \mnemonic{COPI} takes a 16-bit immediate command operand. \mnemonic{COPR} takes any register
	name as its source; its 16-bit value is used as the command operand. A protected-mode \mnemonic{COPR sr} reads
	only the low byte of \reg{sr} and therefore always builds a coprocessor-0 command. Hidden encoded register and
	shift fields are canonicalized by the assembler and are not programmer-visible operands.

	\textbf{Operation:}
	\begin{lstlisting}
if condition_true then
    pending_coprocessor_command = command_operand
endif
\end{lstlisting}

	\textbf{Description:}

	Transfers a command to a coprocessor. For the standard command layout, the high nibble selects the coprocessor and the
	next nibble selects that coprocessor's operation. The remaining low byte is defined by the selected operation.

	\textbf{Write-back:} Writes the command operand to the pending coprocessor-command latch.

	\textbf{Flags:} Preserved by the coprocessor-call instruction.

	\textbf{Exceptions:} May raise privilege-violation faults for supervisor-only operations in protected mode. May raise
	invalid-opcode faults for missing coprocessors and unimplemented coprocessor operations.

	\textbf{Condition field:} If the condition is false, no command is dispatched, the pending command latch is preserved,
	no coprocessor privilege or opcode checks are performed, and status flags are preserved.

	\textbf{Timing:} 6 cycles for the processing-unit call, plus a further 6-cycle coprocessor dispatch slot (12
	cycles in total), plus any bus wait states.

	\textbf{Privilege:} Available in protected mode and supervisor mode. In protected mode, coprocessor-operation nibbles
	0x8--0xF are supervisor-only and raise privilege-violation faults when executed. For \mnemonic{COPR} the privilege
	decision is made on the run-time value of \reg{rS} sampled during decode, not on any encoded field, so whether a
	given \mnemonic{COPR} instruction faults depends on register contents.

	\begin{example}
		\begin{lstlisting}
; Wait for interrupt (Exception Unit)
COPI #0x1900                  ; Coprocessor 1, operation 0x9
; Prefer the WAIT meta-instruction for normal assembly code.

; Return from exception
COPI #0x1A00                  ; Coprocessor 1, operation 0xA
; Prefer the RETE meta-instruction for normal assembly code.

; User exception (system call)
COPI #0x1180                  ; Coprocessor 1, operation 0x1, vector 0x80
; Prefer EXCP #0x80 for normal assembly code.
		\end{lstlisting}
		\caption{Direct Coprocessor Calls}
	\end{example}

\end{instructionbox}

\section{Coprocessor Compatibility}

Standard coprocessor IDs are stable across SIRC-1 models. A CPU model may omit an optional coprocessor, but it must not
assign that coprocessor ID to an incompatible unit. If software calls an absent coprocessor or an operation that the
selected coprocessor does not implement, the CPU raises an invalid-opcode fault.

\begin{table}[H]
	\centering
	\small
	\begin{tabularx}{\textwidth}{llcXX}
		\toprule
		\textbf{ID} & \textbf{Name}        & \textbf{Required} & \textbf{Operations}                                & \textbf{Compatibility rule}                          \\
		\midrule
		0x0         & Processing Unit      & Yes               & Internal                                           & Executes SIRCIS processing-unit instructions         \\
		0x1         & Exception Unit       & Yes               & 0x1, 0x9--0xD (software); 0x0, 0xE--0xF (internal) & Handles reset, faults, interrupts, traps, and links  \\
		0x2         & DMA Unit             & Yes               & 0x8--0xA                                           & Required standard DMA transfer unit                  \\
		0x3         & Integer Maths Unit   & No                & 0x0--0x3                                           & Optional standard integer multiply/divide unit       \\
		0x4--0xF    & Reserved/model space & No                & Model-defined                                      & Reserved unless a CPU model documents the assignment \\
		\bottomrule
	\end{tabularx}
	\caption{Standard Coprocessor IDs}
	\label{tab:standard-coprocessor-ids}
\end{table}

\noindent Exception-unit operation nibble 0x0 is the internal no-operation encoding and operation nibbles 0x2--0x8 are
architecturally undefined. Operations 0xE and 0xF are internal encodings the exception unit uses to dispatch a fault or
a hardware exception; see Chapter~\ref{ch:exceptions} for the internal-instruction mechanism.

\subsection{Model and Revision Policy}

\begin{itemize}
	\item Required coprocessors and documented operations must remain backward compatible within the SIRC-1 family.
	\item Optional standard coprocessors keep their assigned IDs even on models that omit them.
	\item Reserved coprocessor IDs must raise invalid-opcode faults when they are not implemented by the selected CPU model.
	      Reserved operation nibbles and reserved operand values within an implemented coprocessor are architecturally undefined.
	\item Software may probe optional coprocessors by executing the documented operation and handling an invalid-opcode fault, or
	      by using a model-specific capability table if one is provided by system firmware. The invalid-opcode fault records the
	      offending 16-bit coprocessor command word in the \texttt{return\_address} field of the fault-metadata link register
	      (link register \imm{7}), which an emulation handler reads with \mnemonic{ETFR}. The pending command is cleared before
	      the handler runs.
	\item Processing-unit undocumented opcodes are a separate case: they are accepted by the processing unit and have
	      architecturally undefined effects unless a future CPU revision documents them.
\end{itemize}

\section{Exception Unit (Coprocessor 1)}

The Exception Unit handles:
\begin{itemize}
	\item Faults (highest priority)
	\item Hardware exceptions (interrupts, multiple priority levels)
	\item Software exceptions (traps)
	\item Privileged operations
\end{itemize}

\subsection{Common Operations}

\begin{table}[H]
	\centering
	\begin{tabular}{llp{8cm}}
		\toprule
		\textbf{Opcode} & \textbf{Mnemonic} & \textbf{Operation}                        \\
		\midrule
		0x1100-0x11FF   & EXCP              & User exception (trap) at vector 0x60-0xFF \\
		0x1900          & WAIT              & Wait for interrupt                        \\
		0x1A00          & RETE              & Return from exception                     \\
		0x1B00          & RSET              & System reset                              \\
		0x1C00-0x1C37   & ETFR              & Transfer from exception link register     \\
		0x1D00-0x1D37   & ETTR              & Transfer to exception link register       \\
		\bottomrule
	\end{tabular}
	\caption{Exception Unit Operations}
\end{table}

\noindent For \mnemonic{ETFR} and \mnemonic{ETTR}, command bits 7--6 are not decoded, bits 5--4 select the register
(\texttt{00} none, \texttt{01} saved address, \texttt{10} saved status, \texttt{11} both), and bits 3--0 hold the
link-register index, \texttt{0}--\texttt{7}; values with bit 6 or bit 7 set are not distinct operations. Coprocessor-1
operations \texttt{0xE} and \texttt{0xF} are the exception unit's internal encodings for dispatching a fault or a
hardware exception; a software-issued \mnemonic{COPI} or \mnemonic{COPR} naming operation \texttt{0xE} or \texttt{0xF}
is architecturally undefined.

\section{Exception Unit Meta-Instructions}

\begin{instructionbox}{EXCP}{User Exception Meta-Instruction}
	\textbf{Applicability:} SIRC-1, all revisions; requires Coprocessor 1 (Exception Unit), a required coprocessor
	present in all SIRC-1 models.

	\textbf{Assembles to:} \texttt{COPI \#0x1100 | vector}. See Table~\ref{tab:coprocessor-legal-forms} for the full family of legal forms.
	\begin{center}
		\begin{bytefield}[bitwidth=1.1em]{32}
			\bitheader{0-31} \\
			\bitbox{6}{\opcode{0F}} &
			\bitbox{4}{0x0} &
			\bitbox{16}{0x1100 | \#vector} &
			\bitbox{2}{00} &
			\bitbox{4}{Cond}
		\end{bytefield}
	\end{center}

	\textbf{Instruction Fields:}
	\begin{itemize}
		\item \textbf{Register (bits 25--22):} reserved; the assembler emits \texttt{0x0}.
		\item \textbf{Immediate (bits 21--6):} \texttt{0x1100} \texttt{|} \texttt{\#vector}; bits 15--12 select
		      coprocessor 1, bits 11--8 select operation \texttt{0x1}, and bits 7--0 hold \texttt{\#vector},
		      \texttt{0x00}--\texttt{0xFF}.
		\item \textbf{AF (bits 5--4):} fixed at \texttt{00}.
		\item \textbf{Condition field (bits 3--0):} any of the 16 condition-code encodings in
		      Table~\ref{tab:condition-code-encoding}.
	\end{itemize}

	\textbf{Syntax:} \texttt{EXCP \#vector}

	\textbf{Operands:} \texttt{\#vector} is the low 8-bit user exception vector selector.

	\textbf{Operation:}
	\begin{lstlisting}
if condition_true then
    if protected_mode and vector < 0x60 then
        raise privilege_violation_fault
    else if current_exception_level != 0 then
        discard pending_coprocessor_command
    else
        link[0] = { pc, sr, current_exception_level }
        current_exception_level = 1
        pc = vector_table[vector]
    endif
endif
\end{lstlisting}

	\textbf{Description:} Triggers a software exception at the specified user vector, normally 0x60--0xFF. Used by
	programs running in protected mode to invoke system calls or request supervisor-mode services.

	\textbf{Write-back:} Inherits \mnemonic{COPI} command dispatch to the exception unit.

	\textbf{Flags:} Preserved by the coprocessor-call instruction; exception entry effects are defined by the exception
	unit.

	\textbf{Exceptions:} Dispatches a software exception through the exception unit. Raises a privilege-violation fault
	in protected mode if the selected vector is below 0x60. Software exceptions are accepted only at exception level 0;
	an \mnemonic{EXCP} executed inside any exception or fault handler is discarded -- the pending coprocessor command
	is cleared, no vector is fetched, and no fault is raised. Nested system calls from a handler are therefore not
	possible.

	\textbf{Condition field:} If the condition is false, no trap is dispatched and status flags are preserved.

	\textbf{Timing:} 6 cycles for the call plus a 6-cycle exception-unit dispatch slot.

	\textbf{Privilege:} User-callable in protected mode and supervisor mode.

	\begin{example}
		\begin{lstlisting}
LOAD r1, #'A'                 ; Character to print
EXCP #0x80                    ; Trap to OS print routine
		\end{lstlisting}
		\caption{Triggering a Software Exception}
	\end{example}
\end{instructionbox}

\begin{instructionbox}{WAIT}{Wait for Interrupt Meta-Instruction}
	\textbf{Applicability:} SIRC-1, all revisions; requires Coprocessor 1 (Exception Unit), a required coprocessor
	present in all SIRC-1 models.

	\textbf{Assembles to:} \texttt{COPI \#0x1900}. See Table~\ref{tab:coprocessor-legal-forms} for the full family of legal forms.
	\begin{center}
		\begin{bytefield}[bitwidth=1.1em]{32}
			\bitheader{0-31} \\
			\bitbox{6}{\opcode{0F}} &
			\bitbox{4}{0x0} &
			\bitbox{16}{0x1900} &
			\bitbox{2}{00} &
			\bitbox{4}{Cond}
		\end{bytefield}
	\end{center}

	\textbf{Instruction Fields:}
	\begin{itemize}
		\item \textbf{Register (bits 25--22):} reserved; the assembler emits \texttt{0x0}.
		\item \textbf{Immediate (bits 21--6):} fixed at \texttt{0x1900} (coprocessor 1, operation \texttt{0x9},
		      operand byte \texttt{0x00}).
		\item \textbf{AF (bits 5--4):} fixed at \texttt{00}.
		\item \textbf{Condition field (bits 3--0):} any of the 16 condition-code encodings in
		      Table~\ref{tab:condition-code-encoding}.
	\end{itemize}

	\textbf{Syntax:} \texttt{WAIT}

	\textbf{Operands:} None.

	\textbf{Operation:}
	\begin{lstlisting}
if condition_true then
    if protected_mode then
        raise privilege_violation_fault
    else
        stop instruction fetch
        wait until an enabled hardware interrupt line is asserted or the CPU is reset
    endif
endif
\end{lstlisting}

	\textbf{Description:} Stops instruction fetch and places the CPU in a low-power wait state. The CPU issues no bus
	cycles while waiting. It leaves the wait state only when an enabled hardware interrupt line is asserted (NMI is
	always enabled) or when the CPU is reset. Interrupts masked by the hardware interrupt enable bits
	(\texttt{SR.E1}--\texttt{SR.E4}) do not wake the CPU, and faults and software exceptions cannot occur while
	waiting.

	\textbf{Write-back:} Inherits \mnemonic{COPI} command dispatch to the exception unit.

	\textbf{Flags:} Preserved by the coprocessor-call instruction.

	\textbf{Exceptions:} Raises a privilege-violation fault if executed in protected mode.

	\textbf{Condition field:} If the condition is false, the CPU does not enter wait state and status flags are preserved.

	\textbf{Timing:} 6 cycles for the call plus a 6-cycle exception-unit dispatch slot.

	\textbf{Privilege:} Supervisor mode only.

	\begin{example}
		\begin{lstlisting}
idle_loop:
    WAIT                       ; Halt until an interrupt or reset occurs
    BRAN @idle_loop            ; Resume waiting if woken spuriously
		\end{lstlisting}
		\caption{Idling Until an Interrupt}
	\end{example}
\end{instructionbox}

\begin{instructionbox}{RETE}{Return from Exception Meta-Instruction}
	\textbf{Applicability:} SIRC-1, all revisions; requires Coprocessor 1 (Exception Unit), a required coprocessor
	present in all SIRC-1 models.

	\textbf{Assembles to:} \texttt{COPI \#0x1A00}. See Table~\ref{tab:coprocessor-legal-forms} for the full family of legal forms.
	\begin{center}
		\begin{bytefield}[bitwidth=1.1em]{32}
			\bitheader{0-31} \\
			\bitbox{6}{\opcode{0F}} &
			\bitbox{4}{0x0} &
			\bitbox{16}{0x1A00} &
			\bitbox{2}{00} &
			\bitbox{4}{Cond}
		\end{bytefield}
	\end{center}

	\textbf{Instruction Fields:}
	\begin{itemize}
		\item \textbf{Register (bits 25--22):} reserved; the assembler emits \texttt{0x0}.
		\item \textbf{Immediate (bits 21--6):} fixed at \texttt{0x1A00} (coprocessor 1, operation \texttt{0xA},
		      operand byte \texttt{0x00}).
		\item \textbf{AF (bits 5--4):} fixed at \texttt{00}.
		\item \textbf{Condition field (bits 3--0):} any of the 16 condition-code encodings in
		      Table~\ref{tab:condition-code-encoding}.
	\end{itemize}

	\textbf{Syntax:} \texttt{RETE}

	\textbf{Operands:} None.

	\textbf{Operation:}
	\begin{lstlisting}
if condition_true then
    if protected_mode then
        raise privilege_violation_fault
    else
        n = current_exception_level - 1
        pc = link[n].return_address
        sr = link[n].return_status
        current_exception_level = link[n].saved_level
    endif
endif
\end{lstlisting}

	\textbf{Description:} Returns from an exception handler by restoring the saved program counter and status register.
	Executing \mnemonic{RETE} at exception level 0, that is, outside any exception or fault handler, is architecturally
	undefined.

	\textbf{Write-back:} Inherits \mnemonic{COPI} command dispatch to the exception unit.

	\textbf{Flags:} Restored by the exception unit as part of exception return.

	\textbf{Exceptions:} Raises a privilege-violation fault if executed in protected mode.

	\textbf{Condition field:} If the condition is false, no exception return occurs and current status flags are
	preserved.

	\textbf{Timing:} 6 cycles for the call plus a 6-cycle exception-unit dispatch slot.

	\textbf{Privilege:} Supervisor mode only.

	\begin{example}
		\begin{lstlisting}
timer_handler:
    ; ... service the timer interrupt ...
    RETE                       ; Return to the interrupted program
		\end{lstlisting}
		\caption{Returning from an Interrupt Handler}
	\end{example}
\end{instructionbox}

\begin{instructionbox}{RSET}{System Reset Meta-Instruction}
	\textbf{Applicability:} SIRC-1, all revisions; requires Coprocessor 1 (Exception Unit), a required coprocessor
	present in all SIRC-1 models.

	\textbf{Assembles to:} \texttt{COPI \#0x1B00}. See Table~\ref{tab:coprocessor-legal-forms} for the full family of legal forms.
	\begin{center}
		\begin{bytefield}[bitwidth=1.1em]{32}
			\bitheader{0-31} \\
			\bitbox{6}{\opcode{0F}} &
			\bitbox{4}{0x0} &
			\bitbox{16}{0x1B00} &
			\bitbox{2}{00} &
			\bitbox{4}{Cond}
		\end{bytefield}
	\end{center}

	\textbf{Instruction Fields:}
	\begin{itemize}
		\item \textbf{Register (bits 25--22):} reserved; the assembler emits \texttt{0x0}.
		\item \textbf{Immediate (bits 21--6):} fixed at \texttt{0x1B00} (coprocessor 1, operation \texttt{0xB},
		      operand byte \texttt{0x00}).
		\item \textbf{AF (bits 5--4):} fixed at \texttt{00}.
		\item \textbf{Condition field (bits 3--0):} any of the 16 condition-code encodings in
		      Table~\ref{tab:condition-code-encoding}.
	\end{itemize}

	\textbf{Syntax:} \texttt{RSET}

	\textbf{Operands:} None.

	\textbf{Operation:}
	\begin{lstlisting}
if condition_true then
    if protected_mode then
        raise privilege_violation_fault
    else
        assert reset output for 6 cycles
        sr = 0x0000
        pc = vector_table[reset_vector]
    endif
endif
\end{lstlisting}

	\textbf{Description:} Requests a full processor reset. The reset output is asserted and held for 6 cycles, after
	which the exception unit fetches the reset vector, sets \reg{sr} to \texttt{0x0000}, and loads the program counter
	from the vector. General-purpose registers, the \reg{a}, \reg{l}, and \reg{s} pairs, and the exception link
	registers are not affected by reset and hold undefined values after power-on.

	\textbf{Write-back:} Inherits \mnemonic{COPI} command dispatch to the exception unit.

	\textbf{Flags:} Reset by the exception unit as part of reset processing.

	\textbf{Exceptions:} Raises a privilege-violation fault if executed in protected mode.

	\textbf{Condition field:} If the condition is false, no reset occurs and current status flags are preserved.

	\textbf{Timing:} 6 cycles for the call, plus a 6-cycle \texttt{RSTO} hold, plus 6 cycles for the exception unit to
	fetch the reset vector -- 18 cycles from the \mnemonic{RSET} fetch to the first instruction of the reset handler.

	\textbf{Privilege:} Supervisor mode only.

	\begin{example}
		\begin{lstlisting}
; Request a full processor reset from supervisor code
RSET                           ; Reboots through the reset vector
		\end{lstlisting}
		\caption{Requesting a Software Reset}
	\end{example}
\end{instructionbox}

\begin{instructionbox}{ETFR/ETTR}{Exception Link Transfer Meta-Instructions}
	\textbf{Applicability:} SIRC-1, all revisions; requires Coprocessor 1 (Exception Unit), a required coprocessor
	present in all SIRC-1 models.

	\textbf{Assembles to:} \texttt{COPI \#0x1Cxy} or \texttt{COPI \#0x1Dxy}. See Table~\ref{tab:coprocessor-legal-forms} for the full family of legal forms.
	\begin{center}
		\begin{bytefield}[bitwidth=1.1em]{32}
			\bitheader{0-31} \\
			\bitbox{6}{\opcode{0F}} &
			\bitbox{4}{0x0} &
			\bitbox{16}{0x1Cxy} &
			\bitbox{2}{00} &
			\bitbox{4}{Cond}
		\end{bytefield}
	\end{center}
	\begin{center}
		\begin{bytefield}[bitwidth=1.1em]{32}
			\bitheader{0-31} \\
			\bitbox{6}{\opcode{0F}} &
			\bitbox{4}{0x0} &
			\bitbox{16}{0x1Dxy} &
			\bitbox{2}{00} &
			\bitbox{4}{Cond}
		\end{bytefield}
	\end{center}

	\textbf{Instruction Fields:}
	\begin{itemize}
		\item \textbf{Register (bits 25--22):} reserved; the assembler emits \texttt{0x0}.
		\item \textbf{Immediate (bits 21--6):} \texttt{0x1Cxy} (\mnemonic{ETFR}) or \texttt{0x1Dxy} (\mnemonic{ETTR});
		      bits 15--12 select coprocessor 1, bits 11--8 select operation \texttt{0xC} or \texttt{0xD}. Within the
		      low operand byte \texttt{xy}: bits 7--6 are not decoded, bits 5--4 (\texttt{x}) select the register
		      (\texttt{00} none, \texttt{01} saved address only, \texttt{10} saved status only, \texttt{11} both), and
		      bits 3--0 (\texttt{y}) hold \texttt{\#n}, the link-register index \texttt{0}--\texttt{7}.
		\item \textbf{AF (bits 5--4):} fixed at \texttt{00}.
		\item \textbf{Condition field (bits 3--0):} any of the 16 condition-code encodings in
		      Table~\ref{tab:condition-code-encoding}.
	\end{itemize}

	\textbf{Syntax:}
	\begin{lstlisting}
ETFR #n                       ; Transfer saved address and status to a and r7
ETFR a, #n                    ; Transfer saved address to a only
ETFR r7, #n                   ; Transfer saved status to r7 only
ETTR #n                       ; Transfer a and r7 to saved address and status
ETTR #n, a                    ; Transfer a to saved address only
ETTR #n, r7                   ; Transfer r7 to saved status only
\end{lstlisting}

	\textbf{Operands:} \texttt{\#n} selects one of the eight exception link registers and must be in the range
	\imm{0}--\imm{7}. Link register \imm{7} holds fault metadata; link register $L-1$ holds the state saved on entry to
	an exception at level $L$. Values above \imm{7} are architecturally undefined; the assembler currently masks
	\texttt{\#n} to four bits without diagnosing the error. Optional operands select whether to transfer the saved
	return address, saved status register, or both.

	\textbf{Operation:}
	\begin{lstlisting}
if condition_true then
    if protected_mode then
        raise privilege_violation_fault
    else if operation == ETFR then
        if x includes saved address then a = link[n].return_address
        if x includes saved status then r7 = link[n].return_status
    else
        if x includes saved address then link[n].return_address = a
        if x includes saved status then link[n].return_status = r7
    endif
endif
\end{lstlisting}

	\textbf{Description:} \mnemonic{ETFR} copies saved exception state from an exception link register to CPU registers.
	\mnemonic{ETTR} copies CPU register state back to the selected exception link register before returning.

	\textbf{Write-back:} \mnemonic{ETFR} writes selected saved state to \reg{a}, \reg{r7}, or both. \mnemonic{ETTR} writes
	\reg{a}, \reg{r7}, or both to selected saved exception state.

	\textbf{Flags:} Preserved by the coprocessor-call instruction.

	\textbf{Exceptions:} Raises a privilege-violation fault if executed in protected mode.

	\textbf{Condition field:} If the condition is false, no transfer occurs and status flags are preserved.

	\textbf{Timing:} 6 cycles for the call plus a 6-cycle exception-unit dispatch slot.

	\textbf{Privilege:} Supervisor mode only.

	\begin{example}
		\begin{lstlisting}
alignment_fault_handler:
    ETFR a, #6                ; Get faulting return address
    LOAD al, #0x0100          ; Correct the address
    ETTR #6, a                ; Write corrected address back
    RETE
		\end{lstlisting}
		\caption{Correcting a Faulting Return Address}
	\end{example}
\end{instructionbox}

See Chapter~\ref{ch:exceptions} for detailed documentation of exception handling and all exception coprocessor
instructions.

\section{DMA Unit (Coprocessor 2)}
\label{sec:dma-unit}

The required DMA unit provides bulk word transfers using the existing address register pairs and the general-purpose
register file. It is intentionally simple: it has no descriptor table, hidden parameter registers, or autonomous
background execution. When a DMA command is accepted, the processing unit stops fetching further instructions until the
DMA command completes or faults. As of this printing, the reference SIRC-1 emulator has not yet implemented the DMA
unit; software must not assume Coprocessor 2 is present until an emulator release documents it.

DMA operations are supervisor-only because they perform memory access using supervisor-visible address state. In
protected mode, they raise a privilege-violation fault before dispatch.

\subsection{DMA Operations}

\begin{table}[H]
	\centering
	\small
	\begin{tabularx}{\textwidth}{llllX}
		\toprule
		\textbf{Mnemonic} & \textbf{Command} & \textbf{Count}              & \textbf{Registers}                          & \textbf{Operation}                          \\
		\midrule
		\mnemonic{DMAR}   & 0x2800 | operand & \texttt{\#-7}--\texttt{\#7} & \reg{a}/\reg{l}/\reg{s}, \reg{r1}--\reg{r7} & Read memory into registers                  \\
		\mnemonic{DMAW}   & 0x2900 | operand & \texttt{\#-7}--\texttt{\#7} & \reg{a}/\reg{l}/\reg{s}, \reg{r1}--\reg{r7} & Write registers to memory                   \\
		\mnemonic{DMAT}   & 0x2A00 | n       & 0--255                      & \reg{a}, \reg{l}, \reg{r1}--\reg{r7}        & Copy memory at \reg{a} to memory at \reg{l} \\
		\bottomrule
	\end{tabularx}
	\caption{DMA Unit Operations}
	\label{tab:dma-unit-operations}
\end{table}

\noindent For \mnemonic{DMAR} and \mnemonic{DMAW}, the command operand byte is encoded as follows: bits 7--6 select the
address register (\texttt{00} = \reg{a}, \texttt{01} = \reg{l}, \texttt{10} = \reg{s}, \texttt{11} reserved), bit 5
selects direction (\texttt{0} = forward, \texttt{1} = backward), bits 2--0 hold the count magnitude, and bits 4--3 are
reserved and written as zero. Negative counts are encoded as the backward direction bit plus the positive magnitude, not
as a two's-complement count field. A zero count is encoded with the direction bit clear (forward). Note: this DMA
operand field is not the AF address-register field; its \texttt{00}/\texttt{01} assignment differs from
Table~\ref{tab:addr-reg-encoding} in Chapter~\ref{ch:addressing-modes}.

\noindent For all DMA operations, a count of zero is a no-op: no memory access occurs, registers are preserved, and
address registers are not advanced. CPU status flags are preserved. If the condition is false, no DMA command is
dispatched, no coprocessor privilege or opcode checks are performed, and the pending coprocessor command, registers,
memory, address registers, and status flags are preserved; the program counter still advances to the following
instruction.

\noindent DMA timing assumes the command condition is true and every bus access is acknowledged without wait states. The
\mnemonic{COPI} dispatch takes the normal 6 cycles; dispatch to the DMA unit then occupies one further complete 6-cycle
slot beginning at the next phase 0, matching the coprocessor-1 dispatch model (see Table~\ref{tab:coprocessor-common-semantics}).
For a zero count, that slot accepts the command, decodes the count, and resumes instruction fetch, and no further cycles
are needed. For a non-zero count, the same slot performs the first no-wait DMA bus access. Each additional DMA read or
write bus access then takes one further cycle in the no-wait case. Therefore \mnemonic{DMAR} and \mnemonic{DMAW} take
$12 + \max(|n| - 1, 0)$ cycles, and \mnemonic{DMAT} takes $12 + \max(2n - 1, 0)$ cycles. A false condition skips the DMA
command entirely and takes only the normal 6-cycle \mnemonic{COPI} instruction time.

\noindent DMA memory reads and writes can raise data bus and bus-protection faults using DMA read/write bus access
types. If an address-register low word overflows or underflows and \texttt{SR.A} is set, the DMA operation raises a
segment-overflow fault; otherwise the low word wraps modulo $2^{16}$. Partial side effects are possible if a fault
occurs mid-transfer, and software must not assume that re-executing the same DMA instruction restarts the original
operation.

\begin{instructionbox}{DMAR}{DMA Read To Registers Meta-Instruction}
	\textbf{Applicability:} SIRC-1, all revisions; requires Coprocessor 2 (DMA Unit), a required coprocessor present
	in all SIRC-1 models.

	\textbf{Assembles to:} \texttt{COPI \#0x2800 | operand}. See Table~\ref{tab:coprocessor-legal-forms} for the full family of legal forms.
	\begin{center}
		\begin{bytefield}[bitwidth=1.1em]{32}
			\bitheader{0-31} \\
			\bitbox{6}{\opcode{0F}} &
			\bitbox{4}{0x0} &
			\bitbox{16}{0x2800 | operand} &
			\bitbox{2}{00} &
			\bitbox{4}{Cond}
		\end{bytefield}
	\end{center}

	\textbf{Instruction Fields:}
	\begin{itemize}
		\item \textbf{Register (bits 25--22):} reserved; the assembler emits \texttt{0x0}.
		\item \textbf{Immediate (bits 21--6):} \texttt{0x2800} \texttt{|} operand; bits 15--12 select coprocessor 2 and
		      bits 11--8 select operation \texttt{0x8}. Within the low operand byte: bits 7--6 select the address
		      register (\texttt{00} = \reg{a}, \texttt{01} = \reg{l}, \texttt{10} = \reg{s}, \texttt{11} reserved --
		      this is not the AF address-register encoding), bit 5 selects direction (\texttt{0} forward, \texttt{1}
		      backward), bits 4--3 are reserved and written as zero, and bits 2--0 hold the count magnitude,
		      \texttt{0}--\texttt{7}.
		\item \textbf{AF (bits 5--4):} fixed at \texttt{00}.
		\item \textbf{Condition field (bits 3--0):} any of the 16 condition-code encodings in
		      Table~\ref{tab:condition-code-encoding}.
	\end{itemize}

	\textbf{Syntax:} \texttt{DMAR addr, \#n}

	\textbf{Operands:} \texttt{addr} must be \reg{a}, \reg{l}, or \reg{s}. \texttt{\#n} is a literal count operand in the
	range \texttt{\#-7}--\texttt{\#7}. Counts outside that range are invalid DMA operations.

	\textbf{Operation:}
	\begin{lstlisting}
count = abs(n)
if n < 0:
    addr.low = addr.low - count
transfer = addr
for i in 1..count:
    r[i] = memory[transfer]
    transfer.low = transfer.low + 1
if n > 0:
    addr.low = addr.low + count
\end{lstlisting}

	\textbf{Description:} Reads sequential words from the selected address register into \reg{r1} through \reg{rn}.
	Positive counts read from the current address and advance it. Negative counts first move the address register backward
	by the count magnitude, then read the ascending memory block from that new address.

	\textbf{Write-back:} A successful positive-count \mnemonic{DMAR} advances the selected address register by the count
	magnitude. A successful negative-count \mnemonic{DMAR} leaves the selected address register at its decremented start
	address. Address-register high words are not changed by the transfer increments or decrements. Writes \reg{r1}
	through \reg{rn}.

	\textbf{Flags:} Preserved.

	\textbf{Exceptions:} Unsupported DMA operations, reserved operand encodings, and invalid counts raise invalid-opcode
	faults. In protected mode, raises a privilege-violation fault. DMA reads can raise data bus, bus-protection, and
	segment-overflow faults.

	\textbf{Condition field:} If the condition is false, no DMA command is dispatched and registers, memory, address
	registers, and status flags are preserved; the program counter still advances to the following instruction.

	\textbf{Timing:} $12 + \max(|n| - 1, 0)$ cycles when the command condition is true and the DMA read cycles are
	acknowledged without wait states. The $n = 0$ no-op form and one-word transfer forms take 12 cycles. Wait states add
	to the relevant DMA read bus cycle. If the condition is false, the DMA command is not dispatched and the instruction
	takes 6 cycles.

	\textbf{Privilege:} Supervisor mode only.

	\begin{example}
		\begin{lstlisting}
LOAD ah, #0x0040
LOAD al, #0x0000              ; a = 0x00400000, start of source block
DMAR a, #3                    ; r1, r2, r3 = memory[a], memory[a+1], memory[a+2]; a += 3
		\end{lstlisting}
		\caption{Reading a Memory Block into Registers}
	\end{example}
\end{instructionbox}

\begin{instructionbox}{DMAW}{DMA Write From Registers Meta-Instruction}
	\textbf{Applicability:} SIRC-1, all revisions; requires Coprocessor 2 (DMA Unit), a required coprocessor present
	in all SIRC-1 models.

	\textbf{Assembles to:} \texttt{COPI \#0x2900 | operand}. See Table~\ref{tab:coprocessor-legal-forms} for the full family of legal forms.
	\begin{center}
		\begin{bytefield}[bitwidth=1.1em]{32}
			\bitheader{0-31} \\
			\bitbox{6}{\opcode{0F}} &
			\bitbox{4}{0x0} &
			\bitbox{16}{0x2900 | operand} &
			\bitbox{2}{00} &
			\bitbox{4}{Cond}
		\end{bytefield}
	\end{center}

	\textbf{Instruction Fields:}
	\begin{itemize}
		\item \textbf{Register (bits 25--22):} reserved; the assembler emits \texttt{0x0}.
		\item \textbf{Immediate (bits 21--6):} \texttt{0x2900} \texttt{|} operand; bits 15--12 select coprocessor 2 and
		      bits 11--8 select operation \texttt{0x9}. Within the low operand byte: bits 7--6 select the address
		      register (\texttt{00} = \reg{a}, \texttt{01} = \reg{l}, \texttt{10} = \reg{s}, \texttt{11} reserved --
		      this is not the AF address-register encoding), bit 5 selects direction (\texttt{0} forward, \texttt{1}
		      backward), bits 4--3 are reserved and written as zero, and bits 2--0 hold the count magnitude,
		      \texttt{0}--\texttt{7}.
		\item \textbf{AF (bits 5--4):} fixed at \texttt{00}.
		\item \textbf{Condition field (bits 3--0):} any of the 16 condition-code encodings in
		      Table~\ref{tab:condition-code-encoding}.
	\end{itemize}

	\textbf{Syntax:} \texttt{DMAW addr, \#n}

	\textbf{Operands:} \texttt{addr} must be \reg{a}, \reg{l}, or \reg{s}. \texttt{\#n} is a literal count operand in the
	range \texttt{\#-7}--\texttt{\#7}. Counts outside that range are invalid DMA operations.

	\textbf{Operation:}
	\begin{lstlisting}
count = abs(n)
if n < 0:
    addr.low = addr.low - count
transfer = addr
for i in 1..count:
    memory[transfer] = r[i]
    transfer.low = transfer.low + 1
if n > 0:
    addr.low = addr.low + count
\end{lstlisting}

	\textbf{Description:} Writes \reg{r1} through \reg{rn} to sequential words at the selected address register. Positive
	counts write from the current address and advance it. Negative counts first move the address register backward by the
	count magnitude, then write the ascending memory block from that new address. Register order always corresponds to
	ascending memory addresses, so \texttt{DMAW s, \#-n} followed by \texttt{DMAR s, \#n} restores \reg{r1} through
	\reg{rn}.

	\textbf{Write-back:} A successful positive-count \mnemonic{DMAW} advances the selected address register by the count
	magnitude. A successful negative-count \mnemonic{DMAW} leaves the selected address register at its decremented start
	address. Address-register high words are not changed by the transfer increments or decrements. General-purpose
	registers are preserved; \mnemonic{DMAW} only reads them.

	\textbf{Flags:} Preserved.

	\textbf{Exceptions:} Unsupported DMA operations, reserved operand encodings, and invalid counts raise invalid-opcode
	faults. In protected mode, raises a privilege-violation fault. DMA writes can raise data bus, bus-protection, and
	segment-overflow faults.

	\textbf{Condition field:} If the condition is false, no DMA command is dispatched and registers, memory, address
	registers, and status flags are preserved; the program counter still advances to the following instruction.

	\textbf{Timing:} $12 + \max(|n| - 1, 0)$ cycles when the command condition is true and the DMA write cycles are
	acknowledged without wait states. The $n = 0$ no-op form and one-word transfer forms take 12 cycles. Wait states add
	to the relevant DMA write bus cycle. If the condition is false, the DMA command is not dispatched and the instruction
	takes 6 cycles.

	\textbf{Privilege:} Supervisor mode only.

	\begin{example}
		\begin{lstlisting}
LOAD ah, #0x0040
LOAD al, #0x0000              ; a = 0x00400000, start of destination block
DMAW a, #3                    ; memory[a], memory[a+1], memory[a+2] = r1, r2, r3; a += 3
		\end{lstlisting}
		\caption{Writing Registers to a Memory Block}
	\end{example}
\end{instructionbox}

\begin{instructionbox}{DMAT}{DMA Memory Transfer Meta-Instruction}
	\textbf{Applicability:} SIRC-1, all revisions; requires Coprocessor 2 (DMA Unit), a required coprocessor present
	in all SIRC-1 models.

	\textbf{Assembles to:} \texttt{COPI \#0x2A00 | n}. See Table~\ref{tab:coprocessor-legal-forms} for the full family of legal forms.
	\begin{center}
		\begin{bytefield}[bitwidth=1.1em]{32}
			\bitheader{0-31} \\
			\bitbox{6}{\opcode{0F}} &
			\bitbox{4}{0x0} &
			\bitbox{16}{0x2A00 | \#n} &
			\bitbox{2}{00} &
			\bitbox{4}{Cond}
		\end{bytefield}
	\end{center}

	\textbf{Instruction Fields:}
	\begin{itemize}
		\item \textbf{Register (bits 25--22):} reserved; the assembler emits \texttt{0x0}.
		\item \textbf{Immediate (bits 21--6):} \texttt{0x2A00} \texttt{|} \texttt{\#n}; bits 15--12 select coprocessor
		      2, bits 11--8 select operation \texttt{0xA}, and bits 7--0 hold \texttt{\#n}, \texttt{0}--\texttt{255}.
		\item \textbf{AF (bits 5--4):} fixed at \texttt{00}.
		\item \textbf{Condition field (bits 3--0):} any of the 16 condition-code encodings in
		      Table~\ref{tab:condition-code-encoding}.
	\end{itemize}

	\textbf{Syntax:} \texttt{DMAT a, l, \#n}

	\textbf{Operands:} The source and destination address registers must be written explicitly as \reg{a} and \reg{l}.
	\texttt{\#n} is an 8-bit literal count operand in the range 0--255. A count of zero is a no-op.

	\textbf{Operation:}
	\begin{lstlisting}
while remaining > 0:
    chunk = min(remaining, 7)
    read chunk words from memory[a] into r1..r[chunk]
    write chunk words from r1..r[chunk] to memory[l]
    a.low = a.low + chunk
    l.low = l.low + chunk
    remaining = remaining - chunk
\end{lstlisting}

	\textbf{Description:} Copies sequential words from the address in \reg{a} to the address in \reg{l}, using
	\reg{r1}--\reg{r7} as an internal transfer window in chunks of up to seven words.

	\textbf{Write-back:} A successful non-zero \mnemonic{DMAT} advances \reg{a} and \reg{l} by the number of words
	transferred. Address-register high words are not changed by the transfer increments. \reg{r1}--\reg{r7} are
	clobbered; after completion they contain the final transfer window.

	\textbf{Flags:} Preserved.

	\textbf{Exceptions:} Unsupported DMA operations and reserved operand encodings raise invalid-opcode faults. In
	protected mode, raises a privilege-violation fault. DMA reads and writes can raise data bus, bus-protection, and
	segment-overflow faults.

	\textbf{Fault side effects:} Registers, memory, and address registers reflect the last completed word transfer if a
	fault occurs mid-transfer.

	\textbf{Condition field:} If the condition is false, no DMA command is dispatched and registers, memory, address
	registers, and status flags are preserved; the program counter still advances to the following instruction.

	\textbf{Timing:} $12 + \max(2n - 1, 0)$ cycles when the command condition is true and all DMA read and write cycles
	are acknowledged without wait states. The $n = 0$ no-op form takes 12 cycles, and a one-word memory transfer takes 13
	cycles. Wait states add to the relevant DMA read or write bus cycle. Instruction fetch resumes only after the DMA
	operation completes or faults. If the condition is false, the DMA command is not dispatched and the instruction takes 6
	cycles.

	\textbf{Privilege:} Supervisor mode only.

	\textbf{Overlap:} Overlapping source and destination ranges are architecturally undefined unless the source and
	destination addresses are identical.

	\begin{example}
		\begin{lstlisting}
LOAD ah, #0x0040
LOAD al, #0x0000              ; a = 0x00400000, source block
LOAD lh, #0x0050
LOAD ll, #0x0000              ; l = 0x00500000, destination block
DMAT a, l, #16                ; Copy 16 words from memory[a] to memory[l]
		\end{lstlisting}
		\caption{Copying a Memory Block with DMAT}
	\end{example}
\end{instructionbox}

\section{Integer Maths Unit (Coprocessor 3)}
\label{sec:integer-maths-unit}

The optional integer maths unit provides a stable multiply/divide convention. CPU models without the hardware raise an
invalid-opcode fault for these commands, allowing supervisor software to emulate the operation and return as if the
coprocessor call had completed.

Reference source listings for software emulation routines are supplied with the SIRC-1 distribution media. Those
listings are examples of a supervisor compatibility handler, not additional instruction encodings.

\subsection{Maths Operations}

\begin{table}[H]
	\centering
	\small
	\begin{tabularx}{\textwidth}{llllX}
		\toprule
		\textbf{Mnemonic} & \textbf{Command} & \textbf{Inputs}             & \textbf{Outputs}             & \textbf{Operation}              \\
		\midrule
		\mnemonic{MULU}   & 0x3000           & \reg{r1}, \reg{r2}          & \reg{r1}:\reg{r2}, \reg{r3}  & Unsigned 16 x 16 -> 32 multiply \\
		\mnemonic{MULS}   & 0x3100           & \reg{r1}, \reg{r2}          & \reg{r1}:\reg{r2}, \reg{r3}  & Signed 16 x 16 -> 32 multiply   \\
		\mnemonic{DIVU}   & 0x3200           & \reg{r1}:\reg{r2}, \reg{r3} & \reg{r1}, \reg{r2}, \reg{r3} & Unsigned divide                 \\
		\mnemonic{DIVS}   & 0x3300           & \reg{r1}:\reg{r2}, \reg{r3} & \reg{r1}, \reg{r2}, \reg{r3} & Signed divide                   \\
		\bottomrule
	\end{tabularx}
	\caption{Integer Maths Unit Operations}
	\label{tab:integer-maths-unit-operations}
\end{table}

\begin{table}[H]
	\centering
	\begin{tabular}{cll}
		\toprule
		\textbf{Bit} & \textbf{Name} & \textbf{Meaning}             \\
		\midrule
		0            & E             & Any maths error              \\
		1            & DZ            & Divide by zero               \\
		2            & OV            & Quotient overflow            \\
		3--15        & Reserved      & Written as zero, ignore read \\
		\bottomrule
	\end{tabular}
	\caption{Maths Status Word in \reg{r3}}
	\label{tab:maths-status-word}
\end{table}

\noindent Integer maths operations preserve CPU status flags. Missing maths hardware and unsupported maths operations
raise invalid-opcode faults. Divide by zero and quotient overflow are reported in \reg{r3}, not as arithmetic
exceptions. If the condition is false, no maths command is dispatched, no coprocessor opcode check is performed, and
registers and status flags are preserved.

\begin{instructionbox}{MULU/MULS}{Integer Multiply Meta-Instructions}
	\textbf{Applicability:} SIRC-1, all revisions; requires Coprocessor 3 (Integer Maths Unit), an optional
	coprocessor absent on models that do not implement it.

	\textbf{Assembles to:} \texttt{COPI \#0x3000} or \texttt{COPI \#0x3100}. See Table~\ref{tab:coprocessor-legal-forms} for the full family of legal forms.
	\begin{center}
		\begin{bytefield}[bitwidth=1.1em]{32}
			\bitheader{0-31} \\
			\bitbox{6}{\opcode{0F}} &
			\bitbox{4}{0x0} &
			\bitbox{16}{0x3000} &
			\bitbox{2}{00} &
			\bitbox{4}{Cond}
		\end{bytefield}
	\end{center}
	\begin{center}
		\begin{bytefield}[bitwidth=1.1em]{32}
			\bitheader{0-31} \\
			\bitbox{6}{\opcode{0F}} &
			\bitbox{4}{0x0} &
			\bitbox{16}{0x3100} &
			\bitbox{2}{00} &
			\bitbox{4}{Cond}
		\end{bytefield}
	\end{center}

	\textbf{Instruction Fields:}
	\begin{itemize}
		\item \textbf{Register (bits 25--22):} reserved; the assembler emits \texttt{0x0}.
		\item \textbf{Immediate (bits 21--6):} fixed at \texttt{0x3000} (\mnemonic{MULU}) or \texttt{0x3100}
		      (\mnemonic{MULS}); bits 15--12 select coprocessor 3, bits 11--8 select the operation, and bits 7--0 are
		      unused and written as zero.
		\item \textbf{AF (bits 5--4):} fixed at \texttt{00}.
		\item \textbf{Condition field (bits 3--0):} any of the 16 condition-code encodings in
		      Table~\ref{tab:condition-code-encoding}.
	\end{itemize}

	\textbf{Syntax:}
	\begin{lstlisting}
MULU                          ; Unsigned multiply
MULS                          ; Signed multiply
\end{lstlisting}

	\textbf{Operands:} None in assembly syntax. The operations multiply \reg{r1} by \reg{r2}.

	\textbf{Operation:}
	\begin{lstlisting}
MULU:
    product = unsigned(r1) * unsigned(r2)
    r1 = product[31:16]
    r2 = product[15:0]
    r3 = 0

MULS:
    product = signed(r1) * signed(r2)
    r1 = product[31:16]
    r2 = product[15:0]
    r3 = 0
\end{lstlisting}

	\textbf{Description:} Multiplies two 16-bit operands and writes the 32-bit product to \reg{r1}:\reg{r2}, high word
	first. \mnemonic{MULU} treats both inputs as unsigned. \mnemonic{MULS} treats both inputs as two's-complement signed
	values.

	\textbf{Write-back:} Writes \reg{r1}, \reg{r2}, and \reg{r3}.

	\textbf{Flags:} CPU status flags are preserved. \reg{r3} is written as zero.

	\textbf{Exceptions:} Missing maths hardware and unsupported multiply operations raise invalid-opcode faults. These
	operations perform no data-memory bus access.

	\textbf{Condition field:} If the condition is false, no maths command is dispatched and registers and status flags are
	preserved.

	\textbf{Timing:} The \mnemonic{COPI} command dispatch uses the normal 6 execution phases. Hardware maths operations
	may take additional implementation-defined cycles before instruction fetch resumes.

	\textbf{Privilege:} User-callable in protected mode and supervisor mode.

	\begin{example}
		\begin{lstlisting}
LOAD r1, #1000                ; Multiplicand
LOAD r2, #2000                ; Multiplier
MULU                          ; r1:r2 = 1000 * 2000, r3 = 0
		\end{lstlisting}
		\caption{Unsigned 16x16 Multiply}
	\end{example}
\end{instructionbox}

\begin{instructionbox}{DIVU/DIVS}{Integer Divide Meta-Instructions}
	\textbf{Applicability:} SIRC-1, all revisions; requires Coprocessor 3 (Integer Maths Unit), an optional
	coprocessor absent on models that do not implement it.

	\textbf{Assembles to:} \texttt{COPI \#0x3200} or \texttt{COPI \#0x3300}. See Table~\ref{tab:coprocessor-legal-forms} for the full family of legal forms.
	\begin{center}
		\begin{bytefield}[bitwidth=1.1em]{32}
			\bitheader{0-31} \\
			\bitbox{6}{\opcode{0F}} &
			\bitbox{4}{0x0} &
			\bitbox{16}{0x3200} &
			\bitbox{2}{00} &
			\bitbox{4}{Cond}
		\end{bytefield}
	\end{center}
	\begin{center}
		\begin{bytefield}[bitwidth=1.1em]{32}
			\bitheader{0-31} \\
			\bitbox{6}{\opcode{0F}} &
			\bitbox{4}{0x0} &
			\bitbox{16}{0x3300} &
			\bitbox{2}{00} &
			\bitbox{4}{Cond}
		\end{bytefield}
	\end{center}

	\textbf{Instruction Fields:}
	\begin{itemize}
		\item \textbf{Register (bits 25--22):} reserved; the assembler emits \texttt{0x0}.
		\item \textbf{Immediate (bits 21--6):} fixed at \texttt{0x3200} (\mnemonic{DIVU}) or \texttt{0x3300}
		      (\mnemonic{DIVS}); bits 15--12 select coprocessor 3, bits 11--8 select the operation, and bits 7--0 are
		      unused and written as zero.
		\item \textbf{AF (bits 5--4):} fixed at \texttt{00}.
		\item \textbf{Condition field (bits 3--0):} any of the 16 condition-code encodings in
		      Table~\ref{tab:condition-code-encoding}.
	\end{itemize}

	\textbf{Syntax:}
	\begin{lstlisting}
DIVU                          ; Unsigned divide
DIVS                          ; Signed divide
\end{lstlisting}

	\textbf{Operands:} None in assembly syntax. The dividend is \reg{r1}:\reg{r2}, high word first. The divisor is
	\reg{r3}.

	\textbf{Operation:}
	\begin{lstlisting}
if divisor is zero:
    r3 = E | DZ
else if quotient does not fit in 16 bits:
    r3 = E | OV
else:
    r1 = remainder
    r2 = quotient
    r3 = 0
\end{lstlisting}

	\textbf{Description:} Divides a 32-bit dividend by a 16-bit divisor. Successful division writes the remainder to
	\reg{r1}, the quotient to \reg{r2}, and zero to \reg{r3}.

	\textbf{Write-back:} Successful division writes \reg{r1}, \reg{r2}, and \reg{r3}. Divide-by-zero and
	quotient-overflow cases leave \reg{r1}:\reg{r2} unchanged and write the maths status word to \reg{r3}.

	\textbf{Flags:} CPU status flags are preserved. Maths errors are reported through \reg{r3}, not through \reg{sr}.

	\textbf{Exceptions:} Missing maths hardware and unsupported divide operations raise invalid-opcode faults. Divide by
	zero and quotient overflow do not raise arithmetic exceptions. These operations perform no data-memory bus access.

	\textbf{Condition field:} If the condition is false, no maths command is dispatched and registers and status flags are
	preserved.

	\textbf{Timing:} The \mnemonic{COPI} command dispatch uses the normal 6 execution phases. Hardware maths operations
	may take additional implementation-defined cycles before instruction fetch resumes. Software emulation timing is
	defined by the invalid-opcode handler.

	\textbf{Privilege:} User-callable in protected mode and supervisor mode.

	\textbf{Signed division:} \mnemonic{DIVS} treats \reg{r1}:\reg{r2} as a 32-bit two's-complement dividend and
	\reg{r3} as a 16-bit two's-complement divisor. The quotient is truncated toward zero and the remainder has the sign
	of the dividend. The case \texttt{-2147483648 / -1} is quotient overflow.

	\begin{example}
		\begin{lstlisting}
LOAD r1, #0                   ; High word of dividend
LOAD r2, #100                 ; Low word of dividend
LOAD r3, #7                   ; Divisor
DIVU                          ; r1 = remainder, r2 = quotient, r3 = 0 on success
		\end{lstlisting}
		\caption{Unsigned Divide}
	\end{example}
\end{instructionbox}
